The lower and upper bounds now concern the same point moment, but their
variational domains differ. The lower bound uses nonnegative smooth controls;
the upper bound permits signed controls in a completed covariance space.
Identifying the limiting constant requires a proof that these two optimizations
have the same value.

Recall the lower value from Section~\ref{S:unmasked-variational} and the upper
expression from Section~\ref{S:original-signed-score-upper}. We first solve a
finite Hilbert-space optimization problem whose Gram matrix has nonnegative
entries. The argument allows singular matrices and includes the zero control. We
then apply this comparison to the physical means and remove the spatial cutoff.
Theorem~\ref{T:intrinsic-second-moment} states the resulting limit. The
remaining sections derive occupation and spatial-energy bounds for its constant.
For one-dimensional spatial white noise, the spatial energy can be evaluated
explicitly. This gives explicit moment bounds, but does not evaluate the full
moment supremum or identify the auxiliary coefficient-one value with the actual
first-moment constant.
